\section{KuaiRand Logged Benchmark Protocol}\label{app:kuairand}

This protocol instantiates the real-data masking benchmark on top of the
KuaiRand sequential recommendation logs. The goal is to test whether the
Effort-Adjusted Coherence Index ($CI^E$) can surface exposure concentration
and filter-bubble-like drift better than passive detectors and a
smoothed-score baseline.

\subsection*{Data and signals}
The primary state signal is \texttt{watch\_ratio}. Auxiliary feedback signals
include \texttt{click} and \texttt{like}; \texttt{dwell}, \texttt{follow},
and \texttt{share} are reserved for secondary validation when available.
For user $u$ at time $t$, let $P_{u,t}$ denote the inferred preference state,
$A_{u,t}$ the exposure policy, and $\Phi_{u,t}$ the composite feedback signal.
All user windows use \texttt{window\_size = 20}, and video tags are binned from
\texttt{video\_features\_basic\_pure.csv}.
The benchmark keeps only users with at least 20 interactions in the healthy,
bubble, and collapse segments, yielding the 367-user evaluation cohort used in
the main text.

\subsection*{Operational score components}
In KuaiRand, $\sigma_P$ is a normalized deviation of the current watch-ratio
window from the healthy watch-ratio baseline, $\sigma_\Phi$ is a stability term
based on within-window watch-ratio dispersion, and $\sigma_A$ is an operational
surrogate based on normalized exposure-tag entropy. This last term should be
read as a logged-data proxy for action reproducibility, not as a direct
estimator of $e^{-H(\Act\mid\Pot)}$, because the benchmark has no replicated
policy logs.

The proxy hierarchy matters: if direct action logs or replica logs exist, the
Gaussian-style exact interpretation can be approximated; if not, the logged-data
entropy surrogate is the intended operational choice. We keep the report honest
by treating the logged benchmark as evidence for a useful proxy, not as a proof
that the surrogate equals the unobserved conditional entropy. In the ladder used
throughout the paper, KuaiRand is a proxy-level benchmark.

\subsection*{Effort proxy}
The effort proxy is defined as the divergence between the current exposure
distribution and the healthy random-exposure baseline distribution:
\[
  E_{u,t} \coloneqq D_{\mathrm{KL}}
  \left(
    \pi^{\mathrm{current}}_{u,t}\,\middle\|\,
    \pi^{\mathrm{healthy}}_{u,t}
  \right).
\]
This choice is directional, decomposable, and directly interpretable as the
cost of keeping the observed trajectory close to a healthy baseline. In
settings where only aggregate exposure counts are available, total variation
distance may be used as a fallback, but KL is the preferred default.
In the follow-up proxy sweep run for the present revision, the KL proxy also
outperformed the total-variation and exposure-Gini alternatives at the same
healthy-only calibration, so we keep it as the default operational choice here.
We use additive smoothing $10^{-6}$ in the KL computation.

\subsection*{Protocol}
The benchmark is executed in three phases.
\begin{enumerate}[leftmargin=*, label=\textbf{\arabic*.}]
  \item \textbf{Healthy.} Exposure follows the random baseline, diversity is
  unconstrained, and effort remains low.
  \item \textbf{Bubble.} Exposure becomes progressively concentrated, so that
  $E_{u,t}$ rises while standard coherence may remain artificially high.
  \item \textbf{Collapse.} The concentrated exposure regime stops tracking the
  healthy baseline, generic detectors begin to fire, and the effort-corrected
  index should lead the baselines.
\end{enumerate}
\subsection*{Evaluation}
Evaluation is carried out sequentially within each user session and then
batched across users with strict temporal splits. Thresholds for $\Score$,
$\ScoreE$, and EWMA-smoothed $\ScoreE$ are calibrated on healthy windows only,
using the 20th percentile of the healthy score distribution. The reference
effort scale $E_0$ is the median healthy KL divergence. Reported metrics are:
\begin{itemize}[leftmargin=*]
  \item detection rate of the bubble/collapse phases,
  \item median detection delay,
  \item false positives in the healthy phase,
  \item sensitivity to $\lambda$ in $\ScoreE$,
  \item comparison against ADWIN, PageHinkley, KSWIN, NoDrift, and a
    smoothed $CI^E$ baseline (EWMA).
\end{itemize}
The detection metric is user-level rather than AUC-based: a user is counted as
detected for bubble or collapse if the detector issues at least one warning
within a fixed four-window neighborhood before the corresponding annotated phase
boundary. Reported uncertainty intervals are paired bootstrap intervals over
users.
All baseline detectors consume the same effort-adjusted scalar signal,
$1 - \ScoreE$, so the comparison is not biased by detector-specific
inputs. The baseline hyperparameters are fixed a priori: ADWIN $\delta = 0.03$;
PageHinkley $\delta = 0.005$, threshold $= 20$, $\alpha = 0.9999$; KSWIN
window $= 30$, stat $= 10$, $\alpha = 0.001$. When enabled, thresholds are
calibrated on healthy windows only.

\subsection*{Controls}
Required controls include shuffled labels, random phase boundaries, no-
coercion ablations, and removal of auxiliary feedback signals. These controls
ensure that the benchmark does not reduce to a hindsight-tuned logged-bandit
exercise. We additionally report a clipped self-normalized tag-reweighting
sanity check as a weak logged-data control, but we do not present that control
as a causal IPS or counterfactual identification result because the repository
lacks item-level propensities and replica policies.
